\chapter{Idiopathic Intracranial Hypertension}
\index{Idiopathic Intracranial Hypertension}
\label{sec:Idiopathic Intracranial Hypertension}

\section{Overview}
\begin{itemize}[noitemsep]
  \item Idiopathic intracranial hypertension (IIH), formerly called pseudotumor cerebri, is a condition of raised pressure within the skull in the absence of a brain tumor or hydrocephalus
  \item It most commonly affects young, overweight women
  \item It causes headache and visual disturbances; untreated, it can lead to permanent vision loss
\end{itemize}
\section{Causes}
The cause is unknown, but the raised intracranial pressure is thought to result from impaired absorption of cerebrospinal fluid. It is strongly associated with obesity and weight gain. Certain medications (tetracyclines, retinoids such as vitamin A, lithium, growth hormone) and conditions such as polycystic ovary syndrome and venous sinus thrombosis can trigger or mimic it.
\section{Risk Factors}
\begin{itemize}[noitemsep]
  \item Female sex and childbearing age
  \item Obesity, especially recent weight gain
  \item Use of certain medications (tetracycline antibiotics, isotretinoin)
  \item Excess vitamin A intake
  \item Endocrine disorders such as polycystic ovary syndrome
  \item Cerebral venous sinus thrombosis
\end{itemize}
\section{Signs and Symptoms}
\subsection{Common symptoms}
\begin{itemize}[noitemsep]
  \item Headache that is often worse in the morning or with lying down
  \item Pulsatile tinnitus (hearing a whooshing sound in sync with the pulse)
  \item Transient visual obscurations (brief episodes of blurred or grayed vision, often on standing)
  \item Double vision (usually from sixth nerve palsy)
  \item Pulsing visual flashes
\end{itemize}
\subsection{Emergency warning signs}
\begin{itemize}[noitemsep]
  \item Sudden, severe, or progressive vision loss requires urgent evaluation, as it can be permanent
  \item Worsening headache with confusion or vomiting requires emergency assessment
\end{itemize}
\section{Progression and Stages}
Without treatment, intracranial pressure remains elevated and can progressively damage the optic nerves, causing gradual visual field loss and eventually blindness. With treatment and weight loss, the pressure usually normalizes and vision stabilizes or improves. The disease can recur, especially with weight regain.
\section{Complications}
\begin{itemize}[noitemsep]
  \item Optic nerve damage (optic atrophy) and permanent vision loss
  \item Visual field defects and constriction
  \item Chronic headaches and reduced quality of life
  \item Sixth nerve palsy causing double vision
  \item Recurrence of symptoms after treatment
\end{itemize}
\section{Diagnosis}
\begin{itemize}[noitemsep]
  \item Lumbar puncture showing elevated opening pressure with normal cerebrospinal fluid
  \item Ophthalmoscopy showing papilledema (swelling of the optic disc)
  \item Visual field testing and formal perimetry
  \item MRI/MRV of the brain to exclude tumors, hydrocephalus, or venous sinus thrombosis
  \item Optical coherence tomography (OCT) of the optic nerve
  \item Diagnosis requires normal imaging, elevated pressure, and characteristic symptoms
\end{itemize}
\section{Prevention}
\begin{itemize}[noitemsep]
  \item Weight management and prevention of obesity
  \item Cautious use of medications known to trigger IIH
  \item Prompt neurological and ophthalmological evaluation of new headaches with visual symptoms
  \item Regular eye examinations for those with risk factors
\end{itemize}
\section{Treatment Options}
\begin{itemize}[noitemsep]
  \item Weight loss (even 5--10\%) is a cornerstone of management
  \item Carbonic anhydrase inhibitors (acetazolamide) to reduce cerebrospinal fluid production
  \item Diuretics (e.g., furosemide) in some cases
  \item Repeat lumbar punctures for temporary relief
  \item Optic nerve sheath fenestration or cerebrospinal fluid shunt surgery when vision is threatened
  \item Discontinuation of any causative medications
\end{itemize}
\section{Living with It}
\begin{itemize}[noitemsep]
  \item Follow the weight loss plan and treatment regimen
  \item Attend regular ophthalmological monitoring of the optic nerves
  \item Report any new or worsening visual symptoms immediately
  \item Manage headaches with the help of a specialist
  \item Connect with support networks for chronic headache conditions
\end{itemize}
\section{Outlook}
With early diagnosis and treatment, vision loss can usually be prevented and headache improved. Weight loss is associated with the best long-term outcomes. Untreated, the condition carries a real risk of permanent blindness, which is why regular visual monitoring is essential.
